\chapter{The Lefschetz--Verdier trace formula}
\label{ch:trace-formula}

The trace formula is first a theorem with finite coefficients and finite Tor
dimension.  Normalization gives the \(\mathbf Z_\ell\)-version and scalar
extension gives the \(\mathbf Q_\ell\)-version used below.  The finite
Tor-dimension hypothesis is essential at finite level: it makes the relevant
objects dualizable and their categorical traces defined.

\section{Traces of perfect complexes}

For a perfect complex \(C\) over a field and an endomorphism \(u\), set
\[
 \operatorname{Tr}(u\mid C)
 =\sum_i(-1)^i\operatorname{Tr}(H^i(u)\mid H^i(C)).
\]

\begin{lemma}[Perfect-complex supertrace]
\label{lem:perfect-supertrace}
Supertrace is invariant under quasi-isomorphism, additive in distinguished
triangles of endomorphisms, multiplicative under derived tensor product, and
cyclic.  Moreover,
\[
 -\log\det(1-Tu\mid C)
 =\sum_{m\geq1}\operatorname{Tr}(u^m\mid C)\frac{T^m}{m},
\]
where
\[
 \det(1-Tu\mid C)
 =\prod_i\det(1-TH^i(u)\mid H^i(C))^{(-1)^i}.
\]
\end{lemma}

\begin{proof}
Represent \(C\) by a bounded finite-dimensional complex and split it into
cohomology and contractible two-term summands.  Contractible summands have
zero supertrace.  This proves invariance and reduces additivity, cyclicity,
and multiplicativity to the corresponding matrix identities.  After a
splitting-field extension, the last identity is the sum of
\(-\log(1-\alpha T)=\sum_{m\geq1}\alpha^mT^m/m\) over the eigenvalues.
\end{proof}

\section{Cohomological correspondences}

Let
\[
 c=(c_1,c_2):C\longrightarrow X\times X
\]
be a correspondence.  For
\(K\in D^b_{ctf}(X,\Lambda_n)\), a cohomological correspondence is a morphism
\[
 u:c_2^*K\longrightarrow Rc_1^!K.
\]
Here \(D^b_{ctf}\) means bounded constructible cohomology and finite Tor
dimension over \(\Lambda_n\).  If \(c_2\) is proper, the unit for
\((c_2^*,Rc_{2*})\), followed by \(u\) and the counit for
\((Rc_{1!},Rc_1^!)\), gives an endomorphism
\[
 R\Gamma_c(u):R\Gamma_c(X,K)\longrightarrow R\Gamma_c(X,K).
\]
The fixed scheme is
\[
 \operatorname{Fix}(c)
 =C\times_{X\times X,c,\Delta_X}X.
\]
Evaluation against the dualizing complex gives a trace class on this fixed
scheme.  Integration over a proper open-and-closed component
\(\beta\) gives its local term \(\operatorname{LT}_\beta(u)\).

\begin{theorem}[Lefschetz--Verdier trace formula]
\label{thm:lefschetz-verdier}
Let \(c_2\) and \(\operatorname{Fix}(c)\) be proper and let
\(K\in D^b_{ctf}(X,\Lambda_n)\).  For
\(u:c_2^*K\to Rc_1^!K\),
\[
 \operatorname{Tr}(R\Gamma_c(u)\mid R\Gamma_c(X,K))
 =
 \sum_{\beta\in\pi_0(\operatorname{Fix}(c))}
 \operatorname{LT}_\beta(u).
\]
The compatible finite-level identities induce the same formula for
finite-Tor-amplitude \(\mathbf Z_\ell\)-complexes and for constructible
\(\mathbf Q_\ell\)-complexes.
\end{theorem}

\begin{proof}[Finite-level proof capsule]
Compactify the correspondence and extend \(K\) by zero.  Localization removes
boundary pieces with no fixed point.  Deformation of the graph to the normal
cone of the diagonal, followed by purity, specializes the diagonal to the
zero section of its normal bundle.  Homotopy invariance and the Thom trace
reduce the vector-bundle case to the identity correspondence on a point,
where the formula is the ordinary categorical trace.

The reduction uses three separately verified identities for trace classes:
proper pushforward, etale pullback, and additivity in localization triangles.
They follow from the units and counits of the six operations, tensor--Hom
adjunction, and proper base change.  They also prove independence of the
compactification.  This is \cite[Expos\'e III]{SGA5}.

No semisimplicity or division by a group order occurs.  In particular, one
does not resolve a \(\Lambda_n[G]\)-module by a finite complex of permutation
modules; this would fail in general when \(\ell\mid|G|\).  The finite-level
theorem is applied directly in \(D^b_{ctf}(X,\Lambda_n)\).  Reduction
compatibility and normalization then give the adic statement.
\end{proof}

\section{Frobenius fixed points and local terms}

Let \(X_0/\mathbf F_q\) be separated of finite type,
\(X=X_0\times\overline{\mathbf F}_q\), and let \((K,\phi)\) be a Weil
complex.  For \(m\geq1\), use the graph correspondence
\[
 C=X,\qquad (c_1,c_2)=(\operatorname{id}_X,F_X^m).
\]
The second projection \(F_X^m\) is finite, hence proper, and the iterated Weil
map is exactly a cohomological correspondence
\[
 u_m:(F_X^m)^*K\longrightarrow K=R(\operatorname{id}_X)^!K.
\]
Its fixed points are \(X_0(\mathbf F_{q^m})\).

At a fixed point \(y\), write \(F_y^m\) for the endomorphism of \(K_{\bar y}\)
induced by \(u_m\).  This is the global \(F^m\) on the fixed-point stalk.  If
the underlying closed point \(x\) has degree \(d\mid m\), then
\[
 F_y^m=F_x^{\,m/d}.
\]

\begin{lemma}[Frobenius local term]
\label{lem:frobenius-local-term}
Every fixed point is isolated and
\[
 \operatorname{LT}_y(F^m,K)
 =\operatorname{Tr}(F_y^m\mid K_{\bar y})
 :=\sum_j(-1)^j
 \operatorname{Tr}(F_y^m\mid\mathcal H^j(K)_{\bar y}).
\]
\end{lemma}

\begin{proof}
First handle singular \(X_0\).  Topological invariance of the etale topos,
together with functoriality of trace classes under the universal homeomorphism
\((X_0)_{\mathrm{red}}\to X_0\), first reduces us to the reduced scheme.
Choose a finite stratification
\[
 X_0=\coprod_a S_a
\]
by smooth locally closed subschemes defined over \(\mathbf F_q\); this exists
because the field is perfect.  Every \(S_a\) is stable under \(F_X^m\).
Repeated localization for a closed stratum and its open complement, together
with additivity of global traces and trace classes, reduces the assertion to
\((S_a,K|_{S_a})\).  Excision identifies the local term at a fixed point in
\(X\) with the one computed in its smooth stratum.

We may therefore assume \(X\) smooth near \(y\).  Choose regular parameters.
The graph-minus-diagonal map has linear part \(1-dF_X^m\).  Since the
\(q\)-power Frobenius has zero differential, this is the identity, so the
fixed point has local intersection multiplicity one.  In the
deformation-to-the-normal-cone definition of the local term only the
categorical trace on the stalk remains.  Devissage through the cohomology
sheaves gives the alternating sum.
\end{proof}

\begin{theorem}[Grothendieck--Lefschetz trace formula]
\label{thm:GL-trace}
For every \(m\geq1\),
\[
 \boxed{\;
 \sum_{y\in X_0(\mathbf F_{q^m})}
 \operatorname{Tr}(F_y^m\mid K_{\bar y})
 =
 \sum_i(-1)^i
 \operatorname{Tr}(F^m\mid H_c^i(X,K)).
 \;}
\]
If \(y\) has underlying closed point \(x\), its summand uses
\(F_x^{m/\deg x}\), not \(F_x^m\).
\end{theorem}

\begin{proof}
Apply Theorem~\ref{thm:lefschetz-verdier} to
\((\operatorname{id},F_X^m,u_m)\).  The fixed locus is finite.
Lemma~\ref{lem:frobenius-local-term} computes every local term, and the global
endomorphism is the \(m\)-fold Weil action on compactly supported cohomology.
\end{proof}

\begin{corollary}[Point-counting trace formula]
\label{cor:point-count-trace}
For the constant sheaf,
\[
 \#X_0(\mathbf F_{q^m})
 =\sum_i(-1)^i\operatorname{Tr}
 (F^m\mid H_c^i(X,\mathbf Q_\ell)).
\]
\end{corollary}

\begin{proof}
Every stalk is one dimensional and its Frobenius trace is \(1\).
\end{proof}

\section{Closed points and rational points}

For \(x\in|X_0|\), let \(d_x=[k(x):\mathbf F_q]\).

\begin{lemma}[Closed-point/rational-point conversion]
\label{lem:closed-rational-traces}
For every \(m\geq1\),
\[
 \sum_{y\in X_0(\mathbf F_{q^m})}
 \operatorname{Tr}(F_y^m\mid K_{\bar y})
 =
 \sum_{\substack{x\in|X_0|\\d_x\mid m}}
 d_x\operatorname{Tr}(F_x^{m/d_x}\mid K_{\bar x}).
\]
\end{lemma}

\begin{proof}
A degree-\(d\) closed point contributes exactly when \(d\mid m\), and then
gives \(d\) rational geometric points.  Global \(F^m\) goes around their
orbit \(m/d\) times, so its stalk action is \(F_x^{m/d}\).  Sum the \(d\)
equal traces.
\end{proof}

\section{Euler products and determinants}

Define
\[
 L_x(K,T)
 =\prod_j\det(1-F_xT^{d_x}\mid
 \mathcal H^j(K)_{\bar x})^{(-1)^{j+1}},
\qquad
 L(X_0,K,T)=\prod_{x\in|X_0|}L_x(K,T).
\]
The Euler product is coefficientwise meaningful because only points of
bounded degree contribute to any fixed coefficient.

\begin{theorem}[Cohomological expression for \(L\)-functions]
\label{thm:cohomological-L}
For every Weil complex \(K\),
\[
 \boxed{\;
 L(X_0,K,T)
 =
 \prod_i\det(1-TF\mid H_c^i(X,K))^{(-1)^{i+1}}.
 \;}
\]
In particular the Euler product is rational.
\end{theorem}

\begin{proof}
The determinant--trace identity gives
\[
\begin{aligned}
\log L(X_0,K,T)
 &=\sum_x\sum_{r\geq1}
   \operatorname{Tr}(F_x^r\mid K_{\bar x})\frac{T^{rd_x}}r\\
 &=\sum_{m\geq1}\frac{T^m}{m}
   \sum_{d_x\mid m}d_x
   \operatorname{Tr}(F_x^{m/d_x}\mid K_{\bar x})\\
 &=\sum_{m\geq1}\frac{T^m}{m}
   \sum_{y\in X_0(\mathbf F_{q^m})}
   \operatorname{Tr}(F_y^m\mid K_{\bar y})\\
 &=\sum_{m\geq1}\frac{T^m}{m}
   \sum_i(-1)^i\operatorname{Tr}
   (F^m\mid H_c^i(X,K)).
\end{aligned}
\]
The second equality uses \(m=rd_x\), the third is
Lemma~\ref{lem:closed-rational-traces}, and the fourth is
Theorem~\ref{thm:GL-trace}.  Apply Lemma~\ref{lem:perfect-supertrace} again
and exponentiate.  Theorem~\ref{thm:adic-perfectness} makes the
cohomological product finite.
\end{proof}

\begin{corollary}[Cohomological zeta formula]
\label{cor:zeta-cohomological}
\[
 Z(X_0,T)
 :=\exp\left(\sum_{m\geq1}
 \#X_0(\mathbf F_{q^m})\frac{T^m}{m}\right)
 =
 \prod_i\det(1-TF\mid H_c^i(X,\mathbf Q_\ell))^{(-1)^{i+1}}.
\]
Thus \(Z(X_0,T)\in\mathbf Q_\ell(T)\).  Rationality over \(\mathbf Q\),
integrality, and independence of \(\ell\) are proved only after weights.
\end{corollary}

\section{Formalization checklist}

\begin{enumerate}
\item Define traces for dualizable objects and use
\(D^b_{ctf}(X,\Lambda_n)\) at finite level.
\item Package a correspondence as \((C,c_1,c_2,u)\) with
\(u:c_2^*K\to Rc_1^!K\) and \(c_2\) proper.
\item For Frobenius use
\((c_1,c_2)=(\operatorname{id},F_X^m)\).
\item Construct local terms, proper pushforward, localization additivity, and
excision as separate theorems.  Singular schemes are reduced by an
\(\mathbf F_q\)-defined smooth stratification before using regular parameters.
\item Do not use a permutation resolution at finite level.
\item Distinguish global \(F_y^m\) from local \(F_x\), and prove
\(F_y^m=F_x^{m/\deg x}\) before manipulating the Euler product.
\item Derive the \(L\)-function identity by equality of formal logarithms.
\end{enumerate}

